\documentclass[10pt,letterpaper]{article}
% assessment-curricula-common.tex
% Shared preamble for the software-security-assessment-modern-curricula suite.
% EXTRACTION PROVENANCE: the seven uploaded curriculum documents carried
% byte-identical inline preambles except for their \hypersetup{pdftitle=...}
% and \title{...} lines, which stay document-local. The lines below are the
% common lines (originally preamble lines 2-10 and 12-27) copied verbatim,
% in their original order. Documents load this file with
%   % assessment-curricula-common.tex
% Shared preamble for the software-security-assessment-modern-curricula suite.
% EXTRACTION PROVENANCE: the seven uploaded curriculum documents carried
% byte-identical inline preambles except for their \hypersetup{pdftitle=...}
% and \title{...} lines, which stay document-local. The lines below are the
% common lines (originally preamble lines 2-10 and 12-27) copied verbatim,
% in their original order. Documents load this file with
%   % assessment-curricula-common.tex
% Shared preamble for the software-security-assessment-modern-curricula suite.
% EXTRACTION PROVENANCE: the seven uploaded curriculum documents carried
% byte-identical inline preambles except for their \hypersetup{pdftitle=...}
% and \title{...} lines, which stay document-local. The lines below are the
% common lines (originally preamble lines 2-10 and 12-27) copied verbatim,
% in their original order. Documents load this file with
%   \input{assessment-curricula-common.tex}
% immediately after \documentclass, then set their own \hypersetup, \title,
% \author and \date.
\usepackage[T1]{fontenc}
\usepackage[utf8]{inputenc}
\usepackage{lmodern}
\usepackage[letterpaper,margin=0.80in,headheight=18pt,headsep=15pt]{geometry}
\usepackage{microtype,booktabs,longtable,tabularx,array,enumitem,xcolor,listings,tcolorbox,fancyhdr,titlesec,xurl,hyperref}
\definecolor{navy}{HTML}{18324E}
\definecolor{teal}{HTML}{136E77}
\definecolor{pale}{HTML}{EDF5F6}
\definecolor{papergray}{HTML}{F4F6F8}
\titleformat{\section}{\Large\bfseries\color{navy}}{\thesection}{.65em}{}
\titleformat{\subsection}{\large\bfseries\color{teal}}{\thesubsection}{.65em}{}
\setcounter{tocdepth}{2}
\setlength{\parskip}{4pt}
\setlength{\parindent}{0pt}
\emergencystretch=2em
\pagestyle{fancy}
\fancyhf{}
\fancyhead[L]{\footnotesize\color{navy}Software Security Assessment | Curriculum Suite}
\fancyhead[R]{\footnotesize\color{teal}2026 edition}
\fancyfoot[L]{\footnotesize Educational use | Authorized local laboratories}
\fancyfoot[R]{\footnotesize\thepage}
\lstdefinestyle{reviewcode}{basicstyle=\ttfamily\footnotesize,breaklines=true,columns=fullflexible,keepspaces=true,showstringspaces=false,frame=single,rulecolor=\color{teal},backgroundcolor=\color{papergray},aboveskip=8pt,belowskip=8pt}
\lstset{style=reviewcode}
\newtcolorbox{insight}[1][]{colback=pale,colframe=teal,boxrule=.5pt,arc=2mm,#1}
\newtcolorbox{historynote}[1][]{colback=papergray,colframe=navy,boxrule=.5pt,arc=2mm,#1}

% immediately after \documentclass, then set their own \hypersetup, \title,
% \author and \date.
\usepackage[T1]{fontenc}
\usepackage[utf8]{inputenc}
\usepackage{lmodern}
\usepackage[letterpaper,margin=0.80in,headheight=18pt,headsep=15pt]{geometry}
\usepackage{microtype,booktabs,longtable,tabularx,array,enumitem,xcolor,listings,tcolorbox,fancyhdr,titlesec,xurl,hyperref}
\definecolor{navy}{HTML}{18324E}
\definecolor{teal}{HTML}{136E77}
\definecolor{pale}{HTML}{EDF5F6}
\definecolor{papergray}{HTML}{F4F6F8}
\titleformat{\section}{\Large\bfseries\color{navy}}{\thesection}{.65em}{}
\titleformat{\subsection}{\large\bfseries\color{teal}}{\thesubsection}{.65em}{}
\setcounter{tocdepth}{2}
\setlength{\parskip}{4pt}
\setlength{\parindent}{0pt}
\emergencystretch=2em
\pagestyle{fancy}
\fancyhf{}
\fancyhead[L]{\footnotesize\color{navy}Software Security Assessment | Curriculum Suite}
\fancyhead[R]{\footnotesize\color{teal}2026 edition}
\fancyfoot[L]{\footnotesize Educational use | Authorized local laboratories}
\fancyfoot[R]{\footnotesize\thepage}
\lstdefinestyle{reviewcode}{basicstyle=\ttfamily\footnotesize,breaklines=true,columns=fullflexible,keepspaces=true,showstringspaces=false,frame=single,rulecolor=\color{teal},backgroundcolor=\color{papergray},aboveskip=8pt,belowskip=8pt}
\lstset{style=reviewcode}
\newtcolorbox{insight}[1][]{colback=pale,colframe=teal,boxrule=.5pt,arc=2mm,#1}
\newtcolorbox{historynote}[1][]{colback=papergray,colframe=navy,boxrule=.5pt,arc=2mm,#1}

% immediately after \documentclass, then set their own \hypersetup, \title,
% \author and \date.
\usepackage[T1]{fontenc}
\usepackage[utf8]{inputenc}
\usepackage{lmodern}
\usepackage[letterpaper,margin=0.80in,headheight=18pt,headsep=15pt]{geometry}
\usepackage{microtype,booktabs,longtable,tabularx,array,enumitem,xcolor,listings,tcolorbox,fancyhdr,titlesec,xurl,hyperref}
\definecolor{navy}{HTML}{18324E}
\definecolor{teal}{HTML}{136E77}
\definecolor{pale}{HTML}{EDF5F6}
\definecolor{papergray}{HTML}{F4F6F8}
\titleformat{\section}{\Large\bfseries\color{navy}}{\thesection}{.65em}{}
\titleformat{\subsection}{\large\bfseries\color{teal}}{\thesubsection}{.65em}{}
\setcounter{tocdepth}{2}
\setlength{\parskip}{4pt}
\setlength{\parindent}{0pt}
\emergencystretch=2em
\pagestyle{fancy}
\fancyhf{}
\fancyhead[L]{\footnotesize\color{navy}Software Security Assessment | Curriculum Suite}
\fancyhead[R]{\footnotesize\color{teal}2026 edition}
\fancyfoot[L]{\footnotesize Educational use | Authorized local laboratories}
\fancyfoot[R]{\footnotesize\thepage}
\lstdefinestyle{reviewcode}{basicstyle=\ttfamily\footnotesize,breaklines=true,columns=fullflexible,keepspaces=true,showstringspaces=false,frame=single,rulecolor=\color{teal},backgroundcolor=\color{papergray},aboveskip=8pt,belowskip=8pt}
\lstset{style=reviewcode}
\newtcolorbox{insight}[1][]{colback=pale,colframe=teal,boxrule=.5pt,arc=2mm,#1}
\newtcolorbox{historynote}[1][]{colback=papergray,colframe=navy,boxrule=.5pt,arc=2mm,#1}

\hypersetup{colorlinks=true,linkcolor=navy,urlcolor=teal,pdftitle={Software Security Assessment Curriculum Crosswalk}}
\title{\vspace{-1.6cm}\Huge\color{navy}\bfseries Software Security Assessment Curriculum Crosswalk\\[8pt]\large\color{teal}Five-track alignment, prerequisites, source provenance and competency transitions}
\author{Software Security Engineering Curriculum | Independent instructional design}
\date{Version 1.0 | 17 October 2026}
\begin{document}
\maketitle
\begin{center}\small Based on the learning tracks in Dowd, McDonald and Schuh (2007);\
contemporary topics are separately labeled as instructional extensions.\end{center}
\begin{insight}\textbf{Audience and ethics.} This curriculum is for authorized software assessment only. All exercises use instructor-owned source, synthetic input, offline captures, or isolated localhost/VM environments. No third-party probing is required.\end{insight}
\tableofcontents
\clearpage
\section{Source fidelity and interpretation}
The book preface on printed pages xix--xx defines a shared Introduction (Chapters 1--4) and five focused tracks. This crosswalk reproduces those sets without rearranging them, and maps original subject matter to modern educational competencies. It is an independently authored interpretation, not publisher-endorsed courseware.

\begingroup\footnotesize\setlength\tabcolsep{3pt}
\begin{longtable}{>{\raggedright\arraybackslash}p{\dimexpr0.190\textwidth-6pt\relax}>{\raggedright\arraybackslash}p{\dimexpr0.220\textwidth-6pt\relax}>{\raggedright\arraybackslash}p{\dimexpr0.590\textwidth-6pt\relax}}\toprule
\textbf{Path} & \textbf{Textbook chapters, after F1--F4} & \textbf{Primary audit target}\\\midrule\endfirsthead
\toprule \textbf{Path} & \textbf{Textbook chapters, after F1--F4} & \textbf{Primary audit target}\\\midrule\endhead
UNIX & 5, 6, 7, 8, 9, 10, 13 & Audit an intentionally vulnerable local privileged-helper/daemon model; provide a privilege graph, safety patch, and security test evidence.\\\midrule
Windows & 5, 6, 7, 8, 11, 12, 13 & Assess a deliberately vulnerable Windows service implementing a restricted IPC interface and file operations.\\\midrule
Web & 8, 13, 17, 18 & Assess a local multi-tenant REST application with authentication, record access, and deployment settings.\\\midrule
Network Application & 5, 6, 7, 8, 13, 16 & Assess a length-prefixed message protocol and a concurrent local server; defend parser and resource invariants.\\\midrule
Network Analysis & 5, 6, 7, 8, 13, 14, 15, 16 & Assess a local packet parser, stream-state classifier and firewall decision engine using offline captures.\\\midrule
\bottomrule\end{longtable}\endgroup
\section{Foundational competency architecture}
\begingroup\footnotesize\setlength\tabcolsep{3pt}
\begin{longtable}{>{\raggedright\arraybackslash}p{\dimexpr0.080\textwidth-6pt\relax}>{\raggedright\arraybackslash}p{\dimexpr0.260\textwidth-6pt\relax}>{\raggedright\arraybackslash}p{\dimexpr0.400\textwidth-6pt\relax}>{\raggedright\arraybackslash}p{\dimexpr0.260\textwidth-6pt\relax}}\toprule
\textbf{ID} & \textbf{Source} & \textbf{Student capability} & \textbf{Evidence}\\\midrule\endfirsthead
\toprule \textbf{ID} & \textbf{Source} & \textbf{Student capability} & \textbf{Evidence}\\\midrule\endhead
F1 & Ch. 1: Vulnerabilities; Security Policies; Classifying Vulnerabilities & Create an evidence-based classification for six local code excerpts. & classification.csv; 1-page policy assumptions; issue evidence notes\\\midrule
F2 & Ch. 2: Trust Relationships; Threat Modeling; Architecture Modeling & Model a local two-tier service with authentication and storage boundaries. & DFD; threat register; ranked audit entry points\\\midrule
F3 & Ch. 3: Exposure; Insecure Defaults; Protective Measures & Inspect a prepared Compose manifest and reverse-proxy configuration. & exposure matrix; hardened configuration diff; verification log\\\midrule
F4 & Ch. 4: Code-Auditing Strategies; Internal Flow Analysis; Reporting & Trace a supplied input through normalization, permission checks and an unsafe sink. & annotated flow trace; triage record; remediation PR with regression tests\\\midrule
\bottomrule\end{longtable}\endgroup
\section{All-chapter traceability index}
\begingroup\footnotesize\setlength\tabcolsep{3pt}
\begin{longtable}{>{\raggedright\arraybackslash}p{\dimexpr0.080\textwidth-6pt\relax}>{\raggedright\arraybackslash}p{\dimexpr0.410\textwidth-6pt\relax}>{\raggedright\arraybackslash}p{\dimexpr0.300\textwidth-6pt\relax}>{\raggedright\arraybackslash}p{\dimexpr0.210\textwidth-6pt\relax}}\toprule
\textbf{Chapter} & \textbf{Textbook content} & \textbf{Contemporary instructional extension} & \textbf{Track(s)}\\\midrule\endfirsthead
\toprule \textbf{Chapter} & \textbf{Textbook content} & \textbf{Contemporary instructional extension} & \textbf{Track(s)}\\\midrule\endhead
1 & vulnerability classification; security policies; trust relationships; audit vs. black-box testing & CWE and risk-informed prioritization & All tracks\\\midrule
2 & architecture, trust modeling, authentication, authorization and threat identification & STRIDE and cloud-native trust boundaries & All tracks\\\midrule
3 & attack-surface assessment, insecure defaults, secure channels and protective measures & containers, IAM and configuration-as-code & All tracks\\\midrule
4 & scoping, code navigation, candidate points, data flow, desk-checking and reporting & CodeQL, SARIF, SAST triage and CI verification & All tracks\\\midrule
5 & stack/heap overflows, off-by-one errors, mitigations and impact assessment & sanitizers, fuzzing and memory-safe implementation alternatives & UNIX, Windows, Network Application, Network Analysis\\\midrule
6 & integer boundaries, type conversion, promotions, pointers and undefined assumptions & checked arithmetic and contemporary C/C++ toolchains & UNIX, Windows, Network Application, Network Analysis\\\midrule
7 & variable state, control flow, return values, memory allocation and double frees & semantic data flow and ownership/static analysis & UNIX, Windows, Network Application, Network Analysis\\\midrule
8 & strings, NUL, encodings, path/shell metacharacters and injection & contextual encoding, parameterization and Unicode normalization & UNIX, Windows, Web, Network Application, Network Analysis\\\midrule
9 & user/group IDs, privilege dropping, file creation, links and TOCTOU & Linux capabilities, openat2 and namespace boundaries & UNIX\\\midrule
10 & fork/exec, file descriptors, environment arrays, IPC and RPC & systemd hardening, seccomp and container execution boundaries & UNIX\\\midrule
11 & Windows handles, access tokens, security descriptors, filesystem and registry & modern Windows service identities and DLL search safety & Windows\\\midrule
12 & pipes, impersonation, window messaging, RPC and COM/DCOM & named-pipe access checks and Windows RPC hardening & Windows\\\midrule
13 & race conditions, signals, reentrancy, deadlocks and threads & thread sanitizers, asynchronous services and atomicity & UNIX, Windows, Web, Network Application, Network Analysis\\\midrule
14 & IP, UDP, TCP, fragment processing, streams and validation & IPv6, QUIC and protocol differential tests & Network Analysis\\\midrule
15 & stateless/stateful filters, attack surface, fragmentation and spoofing & nftables, stateful inspection and eBPF/XDP & Network Analysis\\\midrule
16 & application parser auditing, HTTP, ISAKMP, ASN.1 and DNS & structured parsing, gRPC, DNS and stateful fuzzing & Network Application, Network Analysis\\\midrule
17 & HTTP, sessions, authentication, authorization, SQLi and XSS & OAuth/OIDC, CSP, BOLA and API authorization & Web\\\midrule
18 & CGI, Perl, PHP, Java, ASP.NET, service-oriented patterns and injection & modern Java, .NET, JS/TS and Python web stacks & Web\\\midrule
\bottomrule\end{longtable}\endgroup
\section{Specialization traceability: source to evaluated evidence}
\subsection{UNIX and Linux Software Security Assessment}
\begingroup\footnotesize\setlength\tabcolsep{3pt}
\begin{longtable}{>{\raggedright\arraybackslash}p{\dimexpr0.220\textwidth-6pt\relax}>{\raggedright\arraybackslash}p{\dimexpr0.200\textwidth-6pt\relax}>{\raggedright\arraybackslash}p{\dimexpr0.370\textwidth-6pt\relax}>{\raggedright\arraybackslash}p{\dimexpr0.210\textwidth-6pt\relax}}\toprule
\textbf{Textbook / original mechanism} & \textbf{Contemporary extension} & \textbf{Competency and lab} & \textbf{Evaluation evidence}\\\midrule\endfirsthead
\toprule \textbf{Textbook / original mechanism} & \textbf{Contemporary extension} & \textbf{Competency and lab} & \textbf{Evaluation evidence}\\\midrule\endhead
Ch. 5: Buffer Overflows; Assessing Memory Corruption Impact & ASan/UBSan, NX, PIE and hardened allocators & U1: Audit length-dependent copy in a local C utility with a 32-byte destination. & annotated C diff; sanitizer transcript; 3-case test table\\\midrule
Ch. 6: Arithmetic Boundary Conditions; Type Conversion Vulnerabilities & C23/C++23-aware checks and integer sanitizer flags & U2: Inspect a signed length that is converted to size\_t before allocation. & conversion lattice; checked-allocation patch; compiler output\\\midrule
Ch. 7: Auditing Functions; Auditing Memory Management & CodeQL data flow and ownership contracts & U3: Review error-path cleanup for allocation, parser failure and early returns. & control-flow graph; allocation ledger; patched implementation\\\midrule
Ch. 8: Shell Metacharacters; Path Metacharacters; NUL Injection & execve argument vectors, argv validation and normalized paths & U4: Replace shell-based file processing with direct argument-vector execution. & sink inventory; fixed code; metacharacter regression suite\\\midrule
Ch. 9: Privilege Model; Dropping Privileges & Linux capabilities and service account isolation & U5: Audit privilege transition ordering in a hypothetical helper process. & privilege-transition state diagram; audit log; test design\\\midrule
Ch. 9: Symbolic Links; Race Conditions; Temporary Files & openat2, O\_NOFOLLOW and dirfd-relative operations & U6: Inspect stat-then-open file access in a controlled directory. & race timeline; revised API calls; safety argument\\\midrule
Ch. 10: Program Invocation; File Descriptors; Environment Arrays & CLOEXEC, posix\_spawn, systemd sandboxing & U7: Audit subprocess creation for inherited secrets and writable handles. & descriptor inventory; subprocess diff; execution test\\\midrule
Ch. 10: Interprocess Communication; Remote Procedure Calls & UNIX-domain peer credentials, namespaces, seccomp & U8: Inspect a local UNIX domain socket with weak directory and socket permissions. & IPC threat model; permission table; test transcript\\\midrule
Ch. 13: Race Conditions; Process Synchronization; Threads & C11 atomics, ThreadSanitizer and lock discipline & U9: Review a non-atomic shared counter controlling a request quota. & interleaving trace; fixed program; TSAN evidence\\\midrule
Ch. 13: Signals; Asynchronous-Safe Code; Threading Vulnerabilities & signal-safe interfaces and systemd sandbox policies & U10: Review a signal handler that writes to a mutable shared logging buffer. & handler review; hardened unit file; regression log\\\midrule
\bottomrule\end{longtable}\endgroup
\subsection{Windows Software Security Assessment}
\begingroup\footnotesize\setlength\tabcolsep{3pt}
\begin{longtable}{>{\raggedright\arraybackslash}p{\dimexpr0.220\textwidth-6pt\relax}>{\raggedright\arraybackslash}p{\dimexpr0.200\textwidth-6pt\relax}>{\raggedright\arraybackslash}p{\dimexpr0.370\textwidth-6pt\relax}>{\raggedright\arraybackslash}p{\dimexpr0.210\textwidth-6pt\relax}}\toprule
\textbf{Textbook / original mechanism} & \textbf{Contemporary extension} & \textbf{Competency and lab} & \textbf{Evaluation evidence}\\\midrule\endfirsthead
\toprule \textbf{Textbook / original mechanism} & \textbf{Contemporary extension} & \textbf{Competency and lab} & \textbf{Evaluation evidence}\\\midrule\endhead
Ch. 5: Buffer Overflows; Protection Mechanisms & Windows ASan, CFG and modern compiler switches & W1: Assess unchecked array copy in a local C++ component. & memory map; corrected code; validation cases\\\midrule
Ch. 6: Signed/Unsigned Conversions; Truncation & checked arithmetic and Windows integer types & W2: Audit DWORD and size\_t conversions near a buffer allocation. & conversion diagram; regression suite; code diff\\\midrule
Ch. 7: Return Value Testing; Allocation Functions & RAII wrappers, SAL annotations and C++ static analysis & W3: Review a routine that opens a HANDLE and exits on multiple branches. & path coverage table; RAII rewrite; static report\\\midrule
Ch. 8: Windows Unicode Functions; Path Metacharacters & UTF-16 and safe process creation & W4: Examine Windows filename normalization and argument quoting in a test helper. & string transformation map; patched helper; tests\\\midrule
Ch. 11: Access Tokens; Security Descriptors & integrity levels, service SIDs and least privilege & W5: Inspect a Windows service object illustrative DACL and token permissions. & ACL matrix; token reasoning; PowerShell verification output\\\midrule
Ch. 11: DLL Loading; File Access; The Registry & safe search directories, ACL inheritance and service settings & W6: Review insecure DLL resolution and overbroad registry-key ACLs. & load trace; ACL before/after; mitigation evidence\\\midrule
Ch. 12: Pipes; Impersonation in Pipes & pipe ACLs and authenticated client security context & W7: Audit a sample service that trusts all named-pipe clients. & pipe trust diagram; fixed checks; access test\\\midrule
Ch. 12: Auditing RPC Applications; COM; DCOM Access Controls & RPC authorization callbacks and COM identity & W8: Trace a privileged service method exposed through RPC/COM in a provided interface definition. & call graph; authorization matrix; test protocol\\\midrule
Ch. 13: Windows Process Synchronization; Threading Vulnerabilities & SRW locks, interlocked operations and concurrency testing & W9: Review two worker threads updating shared authorization state. & interleaving chart; lock protocol; tests\\\midrule
Ch. 13: Reentrancy; Race Conditions; Threading Vulnerabilities & WinDbg, ETW and CodeQL C/C++ & W10: Combine findings from token, IPC and thread reviews in a service assessment. & consolidated report; test runs; residual-risk memo\\\midrule
\bottomrule\end{longtable}\endgroup
\subsection{Modern Web Application Security Assessment}
\begingroup\footnotesize\setlength\tabcolsep{3pt}
\begin{longtable}{>{\raggedright\arraybackslash}p{\dimexpr0.220\textwidth-6pt\relax}>{\raggedright\arraybackslash}p{\dimexpr0.200\textwidth-6pt\relax}>{\raggedright\arraybackslash}p{\dimexpr0.370\textwidth-6pt\relax}>{\raggedright\arraybackslash}p{\dimexpr0.210\textwidth-6pt\relax}}\toprule
\textbf{Textbook / original mechanism} & \textbf{Contemporary extension} & \textbf{Competency and lab} & \textbf{Evaluation evidence}\\\midrule\endfirsthead
\toprule \textbf{Textbook / original mechanism} & \textbf{Contemporary extension} & \textbf{Competency and lab} & \textbf{Evaluation evidence}\\\midrule\endhead
Ch. 8: Metacharacters; SQL Queries; Unicode & parameterized SQL and context-sensitive encoding & WEB1: Trace a request parameter to a SQL query constructed with string concatenation. & source-to-sink trace; SQL patch; tests\\\midrule
Ch. 8: NUL Character Injection; Path Metacharacters & canonical path resolution and safe upload names & WEB2: Review file download endpoint path construction and filename normalization. & path model; corrected handler; regression suite\\\midrule
Ch. 13: Race Conditions; Threading Vulnerabilities & atomic session updates and distributed coordination & WEB3: Audit a coupon-redemption or quota handler susceptible to double-spend race. & state machine; concurrent test; patch evidence\\\midrule
Ch. 13: Synchronization Problems; Race Conditions & session fixation and token lifecycle design & WEB4: Inspect a login and logout workflow with shared session state. & session lifecycle diagram; tests; audit findings\\\midrule
Ch. 17: HTTP; Headers; Methods; Parameters and Forms & RFC 9110 semantics, proxy headers and HTTP/2 translation & WEB5: Audit trust in forwarded client IP and host headers through a reverse proxy. & trust-boundary diagram; server patch; test cases\\\midrule
Ch. 17: Cookies; Sessions; Authorization and Access Control & OIDC, secure cookies and BOLA testing & WEB6: Assess record retrieval that checks login but not record ownership. & access matrix; authorization patch; negative tests\\\midrule
Ch. 17: SQL Injection; Cross-Site Scripting; Auditing Strategy & ASVS 5.0.0 and OWASP Top 10:2025 & WEB7: Audit output encoding for HTML content and validate CSRF controls for state changes. & ASVS mapping; code annotations; test report\\\midrule
Ch. 18: SQL Injection; Shell Invocation; File Access & framework-safe APIs in Python, Java, C\# and Node.js & WEB8: Compare a legacy CGI-style sink with a modern framework handler that reaches the same dangerous operation. & framework comparison; patched handler; regression tests\\\midrule
Ch. 18: REST; Web Application Platforms & GraphQL/gRPC, OAuth/OIDC, API Top 10:2023 & WEB9: Trace an API token from gateway verification to downstream authorization. & token trust map; negative tests; secure API contract\\\midrule
Ch. 18: ASP.NET; Java; PHP; Configuration & CodeQL-supported languages, SARIF and GitHub Actions & WEB10: Create a CI assessment plan using a small modern web example in a CodeQL-supported language. & workflow YAML; SARIF excerpt; disposition register\\\midrule
\bottomrule\end{longtable}\endgroup
\subsection{Network Application Security Assessment}
\begingroup\footnotesize\setlength\tabcolsep{3pt}
\begin{longtable}{>{\raggedright\arraybackslash}p{\dimexpr0.220\textwidth-6pt\relax}>{\raggedright\arraybackslash}p{\dimexpr0.200\textwidth-6pt\relax}>{\raggedright\arraybackslash}p{\dimexpr0.370\textwidth-6pt\relax}>{\raggedright\arraybackslash}p{\dimexpr0.210\textwidth-6pt\relax}}\toprule
\textbf{Textbook / original mechanism} & \textbf{Contemporary extension} & \textbf{Competency and lab} & \textbf{Evaluation evidence}\\\midrule\endfirsthead
\toprule \textbf{Textbook / original mechanism} & \textbf{Contemporary extension} & \textbf{Competency and lab} & \textbf{Evaluation evidence}\\\midrule\endhead
Ch. 5: Buffer Overflows; Assessing Memory Corruption Impact & sanitizers and coverage-guided fuzzing & NA1: Audit a length-prefixed decoder that copies network payload into fixed storage. & wire format; guard patch; sanitizer log\\\midrule
Ch. 5: Heap Overflows; Protection Mechanisms & libFuzzer, AFL++ and fuzz corpus minimization & NA2: Create a harness for a byte-frame parser with malformed lengths. & fuzzer harness; minimized regression corpus; crash report\\\midrule
Ch. 6: Byte Order; Arithmetic Boundaries; Type Conversions & network-byte-order helpers and overflow-safe size math & NA3: Review big-endian field decoding followed by an allocation-size calculation. & parsing table; checked math patch; test transcript\\\midrule
Ch. 7: Auditing Control Flow; Return Value Testing & explicit state machines and property-based testing & NA4: Map a handshake parser with optional fields, failure returns and retries. & state transition table; code patch; coverage map\\\midrule
Ch. 7: Auditing Memory Management; Function Side-Effects & RAII and memory ownership reviews & NA5: Inspect a parser storing pointers into an input buffer after the caller frees it. & ownership graph; revised structures; test cases\\\midrule
Ch. 8: Embedded Delimiters; Metacharacter Filtering & canonical framing and strict UTF-8 & NA6: Audit delimiter parsing when a CRLF sequence occurs in attacker-controlled fields. & parser invariants; negative corpus; diff\\\midrule
Ch. 13: Threads; Race Conditions; Starvation and Deadlocks & async I/O and backpressure & NA7: Audit concurrent request IDs and per-connection quota state. & sequence diagram; synchronization patch; tests\\\midrule
Ch. 13: Reentrancy; Asynchronous-Safe Code & timeouts, cancellation safety and server resource budgets & NA8: Review cleanup after timeouts in a connection worker. & lifecycle table; timeout test; fixed handler\\\midrule
Ch. 16: Auditing Application Protocols; Data Verification & protobuf, gRPC and formal grammar review & NA9: Derive a field-level security contract for a small binary application protocol. & protocol grammar; test oracle; decoder patch\\\midrule
Ch. 16: Header Parsing; ASN.1; Domain Name System & HTTP/3 and DNS parse validation & NA10: Compare two local parsers applied to a benign but ambiguous synthetic message. & differential test; protocol report; negative corpus\\\midrule
\bottomrule\end{longtable}\endgroup
\subsection{Network Analysis Security Assessment}
\begingroup\footnotesize\setlength\tabcolsep{3pt}
\begin{longtable}{>{\raggedright\arraybackslash}p{\dimexpr0.220\textwidth-6pt\relax}>{\raggedright\arraybackslash}p{\dimexpr0.200\textwidth-6pt\relax}>{\raggedright\arraybackslash}p{\dimexpr0.370\textwidth-6pt\relax}>{\raggedright\arraybackslash}p{\dimexpr0.210\textwidth-6pt\relax}}\toprule
\textbf{Textbook / original mechanism} & \textbf{Contemporary extension} & \textbf{Competency and lab} & \textbf{Evaluation evidence}\\\midrule\endfirsthead
\toprule \textbf{Textbook / original mechanism} & \textbf{Contemporary extension} & \textbf{Competency and lab} & \textbf{Evaluation evidence}\\\midrule\endhead
Ch. 5: Buffer Overflows; Heap Overflows & sanitizer-instrumented packet parser & NET1: Inspect a packet-snapshot parser that trusts caplen as on-wire length. & parser diff; sanitizer evidence; packet fixtures\\\midrule
Ch. 6: Byte Order; Arithmetic Boundaries & IPv4/IPv6 header validation & NET2: Audit offset computations for a bounded IPv4 header and optional fields. & offset proof; validated parser; regression tests\\\midrule
Ch. 7: Auditing Control Flow; Function Side-Effects & data-flow proofs for packet classifiers & NET3: Review allow/deny decisions after failed protocol parsing. & decision table; patch; denial tests\\\midrule
Ch. 8: Embedded Delimiters; Character Sets and Unicode & HTTP parser consistency in inspection layers & NET4: Compare raw and normalized target interpretation in a proxy and analyzer. & normalization diagram; tests; parser policy\\\midrule
Ch. 13: Race Conditions; Threading Vulnerabilities & lock sharding and atomic flow state & NET5: Audit expiration and refresh of a shared connection-state entry. & flow-state chart; corrected locking; test log\\\midrule
Ch. 14: IP Packet Structures; IP Options Processing & dual-stack parsing and IPv6 extension headers & NET6: Audit interpretation of packet headers in an offline capture analyzer. & packet parse checklist; sample PCAP; tests\\\midrule
Ch. 14: TCP Connections; TCP Streams; TCP Processing & stateful stream reconstruction and evasion analysis & NET7: Build a trace of TCP sequence overlap in an offline sample. & stream diagram; test corpus; policy specification\\\midrule
Ch. 15: Stateless Firewalls; Stateful Inspection Firewalls & nftables and explicit connection tracking & NET8: Audit a toy firewall rule order and state transition map. & rule matrix; simulator tests; corrective rules\\\midrule
Ch. 16: Auditing Application Protocols; Header Parsing & IDS/IPS inspection consistency and Suricata & NET9: Compare normalized application records with original packet bytes in a local fixture. & inspection path; detection-rule test; findings\\\midrule
Ch. 16: ASN.1; Domain Name System & Zeek, Suricata, eBPF/XDP limitations & NET10: Inspect a bounded DNS message parser and map key fields to monitoring events. & decoder tests; telemetry schema; limits memo\\\midrule
\bottomrule\end{longtable}\endgroup
\section{Competency overlap matrix}
\begingroup\footnotesize\setlength\tabcolsep{3pt}
\begin{longtable}{>{\raggedright\arraybackslash}p{\dimexpr0.390\textwidth-6pt\relax}>{\raggedright\arraybackslash}p{\dimexpr0.120\textwidth-6pt\relax}>{\raggedright\arraybackslash}p{\dimexpr0.120\textwidth-6pt\relax}>{\raggedright\arraybackslash}p{\dimexpr0.120\textwidth-6pt\relax}>{\raggedright\arraybackslash}p{\dimexpr0.120\textwidth-6pt\relax}>{\raggedright\arraybackslash}p{\dimexpr0.130\textwidth-6pt\relax}}\toprule
\textbf{Competency} & \textbf{UNIX} & \textbf{Win} & \textbf{Web} & \textbf{Net app} & \textbf{Net analysis}\\\midrule\endfirsthead
\toprule \textbf{Competency} & \textbf{UNIX} & \textbf{Win} & \textbf{Web} & \textbf{Net app} & \textbf{Net analysis}\\\midrule\endhead
Security policy and threat modeling & Core & Core & Core & Core & Core\\\midrule
Manual data/control-flow auditing & Core & Core & Core & Core & Core\\\midrule
Native memory and integer safety & Core & Core & Bridge & Core & Core\\\midrule
String and interpretation hazards & Core & Core & Core & Core & Core\\\midrule
Concurrency, state and races & Core & Core & Core & Core & Core\\\midrule
POSIX privilege, paths and IPC & Core & Bridge & Bridge & Bridge & Bridge\\\midrule
Windows tokens, DACLs and IPC & Bridge & Core & Bridge & Bridge & Bridge\\\midrule
Web authorization and sessions & Bridge & Bridge & Core & Bridge & Bridge\\\midrule
Protocol parsers, framing and DNS & Bridge & Bridge & Bridge & Core & Core\\\midrule
Packet reassembly and firewall state & Bridge & Bridge & Bridge & Bridge & Core\\\midrule
Source-based remediation \& evidence & Core & Core & Core & Core & Core\\\midrule
\bottomrule\end{longtable}\endgroup
A Bridge entry means a cross-track skill may be introduced for deployment context but is not included in the original preface chapter set. It is not the same as full module coverage.

\section{Topologically ordered prerequisite graph}
The directed dependency graph below is acyclic and includes the original foundational order before specialization. Each specialized chain is compressed here to its chapter-level outcomes; see the week schedule in the matching curriculum for assessment-level dependencies.

\begin{lstlisting}
digraph Curriculum {
  F1 -> F2 -> F3 -> F4;
  F4 -> UNIX5 -> UNIX6 -> UNIX7 -> UNIX8 -> UNIX9
     -> UNIX10 -> UNIX13 -> UNIX_CAP;
  F4 -> WIN5 -> WIN6 -> WIN7 -> WIN8 -> WIN11
     -> WIN12 -> WIN13 -> WIN_CAP;
  F4 -> WEB8 -> WEB13 -> WEB17 -> WEB18 -> WEB_CAP;
  F4 -> NA5 -> NA6 -> NA7 -> NA8 -> NA13 -> NA16 -> NA_CAP;
  F4 -> NET5 -> NET6 -> NET7 -> NET8 -> NET13
     -> NET14 -> NET15 -> NET16 -> NET_CAP;
}
\end{lstlisting}
\subsection{Recommended course-order algorithm}
\begin{itemize}[leftmargin=*,itemsep=2pt]
\item Initialize F1 after introductory programming and basic software architecture assessment.
\item Require F1 -> F2 -> F3 -> F4; do not award track admission on vendor tool familiarity alone.
\item After F4, select exactly one specialization sequence and follow its source-defined chapter order.
\item Require accepted chapter-associated laboratory evidence before each downstream assessment.
\item Complete capstone scoping before remediation; peer-review and regression tests close the final node.
\end{itemize}
\section{Cross-track transition and bridge requirements}
\begingroup\footnotesize\setlength\tabcolsep{3pt}
\begin{longtable}{>{\raggedright\arraybackslash}p{\dimexpr0.220\textwidth-6pt\relax}>{\raggedright\arraybackslash}p{\dimexpr0.330\textwidth-6pt\relax}>{\raggedright\arraybackslash}p{\dimexpr0.450\textwidth-6pt\relax}}\toprule
\textbf{Transition} & \textbf{Competencies retained} & \textbf{Add before advanced practice}\\\midrule\endfirsthead
\toprule \textbf{Transition} & \textbf{Competencies retained} & \textbf{Add before advanced practice}\\\midrule\endhead
UNIX -> Network Application & Native memory safety; C arithmetic; taint and state & Network message framing, DNS/HTTP parsing, protocol fuzzing and remote trust modeling\\\midrule
Windows -> Network Application & C/C++ memory; APIs; concurrency; IPC trust & Wire protocol validation, endian conventions, POSIX comparison and parser fuzzing\\\midrule
Web -> Network Application & Untrusted-input reasoning; request state; auth concepts & C memory models, binary framing, strict protocol parsing, resource bounds\\\midrule
Network Application -> Network Analysis & Parser, concurrency and on-wire application logic & IPv4/IPv6, TCP reassembly, firewall state, packet inspection consistency\\\midrule
UNIX -> Network Analysis & C, process state, concurrency and unsafe input & IP/TCP details, filtering semantics, intrusion inspection and protocol normalization\\\midrule
\bottomrule\end{longtable}\endgroup
\section{Tool-to-competency mapping}
\begingroup\footnotesize\setlength\tabcolsep{3pt}
\begin{longtable}{>{\raggedright\arraybackslash}p{\dimexpr0.260\textwidth-6pt\relax}>{\raggedright\arraybackslash}p{\dimexpr0.360\textwidth-6pt\relax}>{\raggedright\arraybackslash}p{\dimexpr0.380\textwidth-6pt\relax}}\toprule
\textbf{Tool family} & \textbf{What learner must demonstrate} & \textbf{Proof; not simply tool installation}\\\midrule\endfirsthead
\toprule \textbf{Tool family} & \textbf{What learner must demonstrate} & \textbf{Proof; not simply tool installation}\\\midrule\endhead
CodeQL / Semgrep & Explain a tainted flow and correct an actual source weakness & Query reasoning, source path, false-positive analysis, regression tests\\\midrule
AddressSanitizer / UBSan & Identify violated bounds, lifetime or undefined-operation invariants & Minimal reproducer, trace analysis, code patch, clean repeat\\\midrule
Burp Suite / ZAP & Map inputs to app routes and compare authorized states & Authenticated vs unauthorized negative tests on localhost\\\midrule
WinDbg / GDB & Relate runtime state to exact source assumptions & Stack/object state, trace, invariant argument\\\midrule
Wireshark / Zeek / Suricata & Interpret offline message and detection behavior & Packet fixture, parsing record, outcome discrepancies\\\midrule
GitHub Actions / SARIF & Integrate proven assessments with engineering review & Reproducible workflow, result metadata and disposition records\\\midrule
\bottomrule\end{longtable}\endgroup
\section{Historical-to-modern mapping guardrails}
\begingroup\footnotesize\setlength\tabcolsep{3pt}
\begin{longtable}{>{\raggedright\arraybackslash}p{\dimexpr0.230\textwidth-6pt\relax}>{\raggedright\arraybackslash}p{\dimexpr0.310\textwidth-6pt\relax}>{\raggedright\arraybackslash}p{\dimexpr0.460\textwidth-6pt\relax}}\toprule
\textbf{Source era} & \textbf{Retain the principle} & \textbf{Explicit modern interpretation}\\\midrule\endfirsthead
\toprule \textbf{Source era} & \textbf{Retain the principle} & \textbf{Explicit modern interpretation}\\\midrule\endhead
CGI, ASP, Perl and ActiveX & String interpretation, privilege boundary and unsafe sink & Teach as historical mechanisms; contrast currently supported web runtimes\\\midrule
Pre-CFG/NX defenses & Threat depends on architecture and control coverage & Verify modern mitigations but fix root causes regardless\\\midrule
Legacy Windows messaging & Cross-boundary window messages and spoofing & Do not imply old shatter techniques work unchanged on modern Windows\\\midrule
Pre-QUIC HTTP & Parsing ambiguity and request state & Apply RFC 9110 semantics plus transport-specific HTTP/2/3 standards\\\midrule
Classical firewall models & State tracking and parser disagreement & Compare with nftables/eBPF and modern stream inspection\\\midrule
\bottomrule\end{longtable}\endgroup
\section{Assessment, proficiency and assurance}
Assessment weights are consistent across all tracks; successful completion requires at least 70/100 overall, plus no omitted security evidence in the final capstone. Competency mastery is a separate judgment: a passing multiple-choice score alone never certifies independent auditing ability.

\begingroup\footnotesize\setlength\tabcolsep{3pt}
\begin{longtable}{>{\raggedright\arraybackslash}p{\dimexpr0.220\textwidth-6pt\relax}>{\raggedright\arraybackslash}p{\dimexpr0.110\textwidth-6pt\relax}>{\raggedright\arraybackslash}p{\dimexpr0.670\textwidth-6pt\relax}}\toprule
\textbf{Component} & \textbf{Weight} & \textbf{Measurement}\\\midrule\endfirsthead
\toprule \textbf{Component} & \textbf{Weight} & \textbf{Measurement}\\\midrule\endhead
Weekly laboratories & 30\% & Lab notebook, executable negative tests, verifiable artifacts and remediation review\\\midrule
Source-code auditing & 20\% & Source/sink trace, execution conditions, control evaluation and evidence quality\\\midrule
Knowledge assessments & 15\% & Short scenarios, design arguments, false-positive differentiation\\\midrule
Midterm practical & 10\% & Timed code audit in week 8 with independent reproduction\\\midrule
Capstone & 25\% & Scoped independent analysis, corrections, tests and risk communication\\\midrule
\bottomrule\end{longtable}\endgroup
\subsection{Observable proficiency rubric}
\begingroup\footnotesize\setlength\tabcolsep{3pt}
\begin{longtable}{>{\raggedright\arraybackslash}p{\dimexpr0.200\textwidth-6pt\relax}>{\raggedright\arraybackslash}p{\dimexpr0.800\textwidth-6pt\relax}}\toprule
\textbf{Level} & \textbf{Evidence of competency}\\\midrule\endfirsthead
\toprule \textbf{Level} & \textbf{Evidence of competency}\\\midrule\endhead
Foundational & Recognizes a weak boundary with prompting; distinguishes a bug from a security impact.\\\midrule
Intermediate & Independently traces representative input to sensitive operations and implements a validated fix.\\\midrule
Advanced & Audits cross-component flow, considers mitigations and exploits prerequisites, and evaluates false positives.\\\midrule
Professional & Produces reproducible, peer-reviewable evidence; prioritizes risk and demonstrates systemic variant prevention.\\\midrule
\bottomrule\end{longtable}\endgroup
\subsection{Finding classification}
\begingroup\footnotesize\setlength\tabcolsep{3pt}
\begin{longtable}{>{\raggedright\arraybackslash}p{\dimexpr0.220\textwidth-6pt\relax}>{\raggedright\arraybackslash}p{\dimexpr0.780\textwidth-6pt\relax}}\toprule
\textbf{Status} & \textbf{Acceptance criteria}\\\midrule\endfirsthead
\toprule \textbf{Status} & \textbf{Acceptance criteria}\\\midrule\endhead
Suspected & Potential weakness with insufficient reachability or security impact evidence; explicitly mark as unconfirmed.\\\midrule
Confirmed & Specific trust violation, reachable attacker conditions and controlled validation or logically sufficient code proof.\\\midrule
False positive & Demonstrated non-reachability or valid defense, supported with executable evidence and control-flow reasoning.\\\midrule
Mitigated & Underlying code issue may remain but documented compensating control blocks the defined abuse path; residual risk required.\\\midrule
Fixed & Root cause is patched, adversarial regression tests pass, and the affected path is analyzed after the change.\\\midrule
\bottomrule\end{longtable}\endgroup
Severity is not the same as exploitability or priority: CVSS characterizes severity; EPSS estimates exploitation likelihood for scored CVEs; impact, exposure, business context, and confidence still require expert evaluation. Do not generate fabricated EPSS values for hypothetical educational defects.

\section{References and verified technical resources}
Status review: 17 October 2026. Stable OWASP ASVS is 5.0.0; the 2025 OWASP Top 10 is released. NIST SP 800-218/SSDF 1.1 is final, while SP 800-218 Rev. 1 / SSDF 1.2 remained listed as draft at review time. Verify moving tool compatibility before deploying laboratories.

\noindent\textbf{Dowd, McDonald, and Schuh. The Art of Software Security Assessment (2007).} Textbook; preface and Chapters 1--18.\par\small\url{https://www.oreilly.com/library/view/the-art-of/0321444426/}\par\normalsize\medskip
\noindent\textbf{MITRE. CWE: Common Weakness Enumeration.} Weakness descriptions and relationships.\par\small\url{https://cwe.mitre.org/}\par\normalsize\medskip
\noindent\textbf{NIST. SP 800-218 (SSDF 1.1, final, 2022).} Secure software development activities.\par\small\url{https://csrc.nist.gov/pubs/sp/800/218/final}\par\normalsize\medskip
\noindent\textbf{NIST. SP 800-218 Rev. 1 (SSDF 1.2, draft as of 2026-10).} Research-only draft; not a replacement for final 1.1.\par\small\url{https://csrc.nist.gov/projects/ssdf/publications}\par\normalsize\medskip
\noindent\textbf{OWASP. ASVS 5.0.0 (stable).} Verifiable application security requirements.\par\small\url{https://owasp.org/www-project-application-security-verification-standard/}\par\normalsize\medskip
\noindent\textbf{OWASP. Top 10:2025 (released).} Web risk awareness categories, not a testing checklist.\par\small\url{https://top10.owasp.org/2025/}\par\normalsize\medskip
\noindent\textbf{OWASP. Web Security Testing Guide (Stable).} Web testing methodologies.\par\small\url{https://wstg.owasp.org/}\par\normalsize\medskip
\noindent\textbf{OWASP. API Security Top 10:2023.} API risk awareness.\par\small\url{https://owasp.org/API-Security/editions/2023/en/0x00-header/}\par\normalsize\medskip
\noindent\textbf{GitHub. CodeQL supported languages.} Language coverage and tool limitations.\par\small\url{https://codeql.github.com/docs/codeql-overview/supported-languages-and-frameworks/}\par\normalsize\medskip
\noindent\textbf{GitHub. CodeQL compiled-language build modes.} CodeQL database production.\par\small\url{https://docs.github.com/en/code-security/concepts/code-scanning/codeql/codeql-for-compiled-languages}\par\normalsize\medskip
\noindent\textbf{FIRST. CVSS 4.0 specification.} Severity characterization, not a probability of exploitation.\par\small\url{https://www.first.org/cvss/v4.0/specification-document}\par\normalsize\medskip
\noindent\textbf{FIRST. Exploit Prediction Scoring System.} Exploit prediction and prioritization.\par\small\url{https://www.first.org/epss/}\par\normalsize\medskip
\noindent\textbf{IETF. RFC 9110 HTTP Semantics.} HTTP behavior shared across versions.\par\small\url{https://www.rfc-editor.org/rfc/rfc9110}\par\normalsize\medskip
\noindent\textbf{IETF. RFC 9114 HTTP/3.} HTTP/3 mapping over QUIC.\par\small\url{https://www.rfc-editor.org/rfc/rfc9114}\par\normalsize\medskip
\noindent\textbf{IETF. RFC 9000 QUIC Transport.} QUIC transport behavior.\par\small\url{https://www.rfc-editor.org/rfc/rfc9000}\par\normalsize\medskip
\noindent\textbf{GitHub. SARIF support for code scanning.} Structured static analysis evidence.\par\small\url{https://docs.github.com/en/code-security/code-scanning/integrating-with-code-scanning/sarif-support-for-code-scanning}\par\normalsize\medskip
\noindent\textbf{Linux. openat2 manual page.} Directory-relative path resolution controls.\par\small\url{https://man7.org/linux/man-pages/man2/openat2.2.html}\par\normalsize\medskip
\noindent\textbf{Microsoft. Windows security documentation.} Security descriptors, access tokens, RPC and ACLs.\par\small\url{https://learn.microsoft.com/en-us/windows/security/}\par\normalsize\medskip
\noindent\textbf{LLVM. AddressSanitizer documentation.} Memory error checking.\par\small\url{https://clang.llvm.org/docs/AddressSanitizer.html}\par\normalsize\medskip
\noindent\textbf{SEI CERT. Secure Coding Standards.} Implementation standards.\par\small\url{https://wiki.sei.cmu.edu/confluence/display/seccode}\par\normalsize\medskip
\end{document}
